% Luc Destombe --- licence CC BY-NC-SA 4.0
% https://creativecommons.org/licenses/by-nc-sa/4.0/deed.fr
% Fichier autonome : compiler avec pdflatex, deux fois (signets, nombre de pages).
\documentclass[a4paper,10pt]{article}
\makeatletter
\newif\iflycee@niveaux
\newif\iflycee@palatino
\lycee@palatinotrue

\RequirePackage[utf8]{inputenc}
\RequirePackage[T1]{fontenc}
\RequirePackage[french]{babel}
\RequirePackage{etoolbox}

\newcommand{\ShorthandsOff}{\@ifundefined{shorthandoff}{}{\shorthandoff{;:!?}}}
\newcommand{\ShorthandsOn}{\@ifundefined{shorthandon}{}{\shorthandon{;:!?}}}
\ShorthandsOff
\AtBeginDocument{\ShorthandsOn}
\AtBeginEnvironment{tikzpicture}{\ShorthandsOff}

\RequirePackage{geometry}
\RequirePackage{fancyhdr}
\RequirePackage{lastpage}
\RequirePackage{multicol}
\RequirePackage{array}
\RequirePackage{enumitem}
\RequirePackage{xcolor}
\RequirePackage{needspace}

\RequirePackage{amsmath,amssymb}
\RequirePackage{mathpazo}
\DeclareMathSizes{10.95}{10.95}{8}{5.5}
\begingroup\catcode`\_=\active \catcode`\^=\active
\gdef_#1{\sb{\scriptscriptstyle #1}}
\gdef^#1{\sp{\scriptscriptstyle\lycee@moinsepais #1}}
\endgroup
\newcommand\lycee@moins{\mathbin{%
  \setbox\z@\hbox{$\m@th\scriptscriptstyle\lycee@moinsorig$}%
  \dimen@=.55\wd\z@ \advance\dimen@-.5pt
  \hbox to.75\wd\z@{\hss$\m@th\scriptscriptstyle\vcenter{\hbox{%
    \tikz\draw[line width=.5pt,line cap=round](0,0)--(\the\dimen@,0);}}$\hss}}}
\begingroup\lccode`\~=`\- \lowercase{\endgroup\let~}\lycee@moins
\newcommand\lycee@moinsepais{\mathcode`\-="8000 }
\AtBeginDocument{\mathchardef\lycee@moinsorig=\mathcode`\-\relax}
\AtBeginDocument{\catcode`\_=12 \mathcode`\_="8000
  \catcode`\^=12 \mathcode`\^="8000 }
\newcommand\lycee@exposants{%
  \fontdimen13\textfont2=.48\fontdimen6\textfont2
  \fontdimen14\textfont2=.48\fontdimen6\textfont2
  \fontdimen15\textfont2=.48\fontdimen6\textfont2
  \fontdimen13\scriptfont2=.48\fontdimen6\scriptfont2
  \fontdimen14\scriptfont2=.48\fontdimen6\scriptfont2
  \fontdimen15\scriptfont2=.48\fontdimen6\scriptfont2}
\every@math@size\expandafter{\the\every@math@size\lycee@exposants}
\newlength{\retraitcalculs}
\setlength{\retraitcalculs}{2em}

\RequirePackage{tikz}
\usetikzlibrary{arrows.meta,calc,babel,shapes.misc}
\RequirePackage[most]{tcolorbox}
\tcbuselibrary{skins,breakable}

\definecolor{coulDef}{HTML}{1F4E79}
\definecolor{coulProp}{HTML}{7030A0}
\definecolor{coulRem}{HTML}{C55A11}
\definecolor{coulEx}{HTML}{2E7D32}
\definecolor{coulExo}{HTML}{B03A2E}
\definecolor{coulPart}{HTML}{1F4E79}
\definecolor{coulMeth}{HTML}{1B6E4F}
\definecolor{coulCourbe}{HTML}{0B5ED7}
\definecolor{coulCourbeB}{HTML}{C00000}
\definecolor{coulCourbeC}{HTML}{2E7D32}
\definecolor{coulGrille}{HTML}{8C8C8C}
\definecolor{coulRep}{HTML}{1B6E4F}
\definecolor{coulBar}{HTML}{2E7D32}

\newcommand{\R}{\mathbb{R}}
\newcommand{\Rs}{\mathbb{R}^{*}}
\renewcommand{\le}{\leqslant}
\renewcommand{\ge}{\geqslant}
\newcommand{\vect}[1]{\overrightarrow{#1}}
\newcommand{\norme}[1]{\left\lVert #1 \right\rVert}
\newcommand{\intervalle}[2]{\left[#1\,;#2\right]}
\RequirePackage{cancel}
\renewcommand{\CancelColor}{\color{coulExo}}
\newcommand{\fois}[1]{\times{\color{coulExo}#1}}
\newcommand{\barre}[1]{\cancel{#1}}

\tikzset{grille/.style={coulGrille,line width=0.4pt}}
\newcommand{\quadrillage}[4]{\draw[grille,step=1] (#1,#2) grid (#3,#4);}
\tikzset{
  croix/.style={cross out,draw,line width=0.6pt,inner sep=0pt,minimum size=4pt},
  croixdroite/.style={croix,rotate=45,line width=0.8pt},
}
\newcommand{\pt}[3]{\path (#1) node[croix]{} node[#2,font=\small,inner sep=2.5pt]{$#3$};}

\newcommand{\lycee@repere}[4]{%
  \draw[grille,step=1] (#1,#2) grid (#3,#4);
  \draw[-{Stealth[length=2mm]},thick] (#1,0)--({#3+0.4},0) node[below left=-1pt]{$x$};
  \draw[-{Stealth[length=2mm]},thick] (0,#2)--(0,{#4+0.4}) node[below left=-1pt]{$y$};
  \node[below left,font=\scriptsize,inner sep=1.5pt] at (0,0) {$0$};}
\newcommand{\lycee@gradx}[1]{\foreach \k/\l in {#1}{%
  \draw (\k,0.07)--(\k,-0.07) node[below,font=\scriptsize,inner sep=1.5pt]{$\l$};}}
\newcommand{\lycee@grady}[1]{\foreach \k/\l in {#1}{%
  \draw (0.07,\k)--(-0.07,\k) node[left,font=\scriptsize,inner sep=1.5pt]{$\l$};}}
\newcommand{\lycee@intff}[2]{\left[#1\,;#2\right]}
\newcommand{\lycee@intfo}[2]{\left[#1\,;#2\right[}
\newcommand{\lycee@intof}[2]{\left]#1\,;#2\right]}
\newcommand{\lycee@intoo}[2]{\left]#1\,;#2\right[}
\newcommand{\lycee@ens}[1]{\left\{#1\right\}}
\AtBeginDocument{%
  \providecommand{\repere}{\lycee@repere}%
  \providecommand{\gradx}{\lycee@gradx}%
  \providecommand{\grady}{\lycee@grady}%
  \providecommand{\Cf}{\mathcal{C}\sb{\scriptscriptstyle f}}%
  \providecommand{\Cg}{\mathcal{C}\sb{\scriptscriptstyle g}}%
  \providecommand{\intff}{\lycee@intff}%
  \providecommand{\intfo}{\lycee@intfo}%
  \providecommand{\intof}{\lycee@intof}%
  \providecommand{\intoo}{\lycee@intoo}%
  \providecommand{\ens}{\lycee@ens}}

\newlist{questions}{enumerate}{3}
\setlist[questions,1]{label=\arabic*\textdegree),leftmargin=*,itemsep=2pt,topsep=3pt}
\setlist[questions,2]{label=\alph*),leftmargin=*,itemsep=1pt,topsep=2pt}
\setlist[questions,3]{label=\roman*),leftmargin=*,itemsep=1pt,topsep=1pt}
\setlist[itemize]{label=\textbullet,leftmargin=*,itemsep=2pt,topsep=3pt}

\newcommand{\besoinplace}[1]{\par\penalty\@M\begingroup
  \dimen@=#1\relax\dimen@ii=\pagegoal\advance\dimen@ii-\pagetotal
  \ifdim\dimen@>\dimen@ii\ifdim\dimen@ii>\z@\vfil\fi\break\fi\endgroup}

\newcommand{\lycee@pied}{\thepage/\pageref{LastPage}}
\newcommand{\entete}[2]{%
  \pagestyle{fancy}\fancyhf{}%
  \fancyhead[L]{\footnotesize\color{coulPart}\textbf{#1}}%
  \fancyhead[R]{\footnotesize\color{coulPart}\textbf{#2}}%
  \fancyfoot[C]{\footnotesize\lycee@pied}%
  \renewcommand{\headrulewidth}{0.4pt}%
  \renewcommand{\headrule}{\hbox to\headwidth{%
    \color{coulPart!40}\leaders\hrule height \headrulewidth\hfill}}}

\RequirePackage{ccicons}
\newcommand{\lycee@licence}{{\scriptsize\color{gray}%
  \copyright\ Luc Destombe --- \ccbyncsa\ CC BY-NC-SA 4.0}}
\newif\iflycee@licence \lycee@licencetrue
\newcommand{\sanslicence}{\lycee@licencefalse}
\AtBeginDocument{\iflycee@licence\AddToHookNext{shipout/foreground}{%
  \put(0,0){\rlap{\kern\dimexpr1in+\hoffset+\oddsidemargin+\textwidth\relax
    \llap{\raisebox{-\dimexpr1in+\voffset+\topmargin+\headheight+\headsep
      +\textheight+\footskip\relax}{\lycee@licence}}}}}\fi}

\newcommand{\formulesserrees}{%
  \setlength{\abovedisplayskip}{3pt}\setlength{\belowdisplayskip}{3pt}%
  \setlength{\abovedisplayshortskip}{2pt}\setlength{\belowdisplayshortskip}{3pt}}

\setlength{\parindent}{0pt}

\RequirePackage[implicit=false,hidelinks,pdfencoding=auto,psdextra,bookmarksopen,
  bookmarksopenlevel=1,pdfcreator={},pdfproducer={}]{hyperref}
\RequirePackage{bookmark}
\hypersetup{pdfauthor={Luc Destombe},pdfkeywords={CC BY-NC-SA 4.0}}
\newcounter{lycee@signet}
\newcommand{\signet}[3][]{\stepcounter{lycee@signet}%
  \edef\lycee@dest{\if\relax\detokenize{#1}\relax
    signet.\arabic{lycee@signet}\else#1\fi}%
  \hypertarget{\lycee@dest}{}\bookmark[level=#2,dest=\lycee@dest]{#3}}
\pdfstringdefDisableCommands{%
  \def\R{R}\def\vect#1{#1}\def\Cf{Cf}\def\Cg{Cg}\def\fois#1{×#1}%
  \def\barre#1{#1}\def\intervalle#1#2{[#1;#2]}\def\intff#1#2{[#1;#2]}%
  \def\textsuperscript#1{#1}\def\dfrac#1#2{#1/#2}\def\frac#1#2{#1/#2}%
  \def\vec#1{#1⃗}}%
\RequirePackage[official]{eurosym}

\geometry{top=0.8cm,bottom=1.3cm,left=1cm,right=1cm,footskip=0.6cm}

\pagestyle{fancy}
\fancyhf{}
\renewcommand{\headrulewidth}{0pt}
\fancyfoot[C]{\footnotesize Page \thepage{} / \pageref{LastPage}}

\newcommand{\titrefiche}[2]{%
  \begin{center}
    \Large\textbf{#1}\\[2pt]
    \large\textbf{#2}\\[2pt]
  \end{center}
  \vspace{0.10cm}}

\setlength{\columnseprule}{0.4pt}
\setlength{\columnsep}{15pt}
\renewcommand{\columnseprulecolor}{\color{black!45}}
\setlength{\multicolsep}{2pt}

\newcounter{partiefiche}
\renewcommand{\thepartiefiche}{\Roman{partiefiche}}
\newif\ifapresbandeau
\newcommand{\lycee@bandeau}[1]{%
  \par\penalty-600
  \vspace{0.08cm}%
  \noindent\colorbox{coulPart}{%
    \parbox{\dimexpr\linewidth-2\fboxsep\relax}{%
      \color{white}\bfseries\raggedright\strut#1\strut}}%
  \nopagebreak\par\nopagebreak
  \vspace{0.06cm}\nopagebreak
  \apresbandeautrue}
\newcommand{\partiefiche}[1]{%
  \refstepcounter{partiefiche}%
  \lycee@bandeau{\signet{1}{\thepartiefiche{} --- #1}\thepartiefiche{} --- #1}}
\newcommand{\partiefichesansnum}[1]{\lycee@bandeau{\signet{1}{#1}#1}}

\newcounter{exo}
\newlength{\espaceexo}\setlength{\espaceexo}{7pt}
\newlength{\espacetitre}\setlength{\espacetitre}{2pt}
\newcommand{\exo}[1][\relax]{%
  \par
  \ifapresbandeau\apresbandeaufalse\else\penalty-150\addvspace{\espaceexo}\fi
  \refstepcounter{exo}%
  \noindent\signet[exercice-\arabic{exo}]{2}{Exercice \arabic{exo}}%
  \textbf{\color{coulExo}\underline{Exercice \theexo}}%
  \ifx\relax#1\relax\else\textbf{ : #1}\fi
  \par\nopagebreak\vspace{\espacetitre}\nopagebreak
  \ignorespaces}

\setlist[enumerate]{nosep,leftmargin=*,itemsep=2pt,topsep=2pt}
\setlist[itemize]{nosep,leftmargin=*,topsep=2pt}
\newcommand{\choix}[3]{\par\hspace*{1em}%
  \makebox[0.3\linewidth][l]{a) #1}\makebox[0.3\linewidth][l]{b) #2}c) #3}
\newcommand{\deux}[2]{\par\hspace*{0.6em}%
  \makebox[0.48\linewidth][l]{#1}#2}
\newcommand{\trois}[3]{\par\hspace*{0.6em}%
  \makebox[0.32\linewidth][l]{#1}\makebox[0.32\linewidth][l]{#2}#3}

\renewenvironment{center}{\par\addvspace{3pt}\centering}{\par\addvspace{3pt}}
\formulesserrees
\AtBeginDocument{\formulesserrees}
\clubpenalty=10000
\widowpenalty=10000
\displaywidowpenalty=10000
\brokenpenalty=10000
\setlength{\parskip}{0pt}

\RequirePackage{ragged2e}
\setlength{\RaggedRightRightskip}{0pt plus 1fil}
\AtBeginDocument{\RaggedRight\binoppenalty=10000 \relpenalty=10000 }
\g@addto@macro\@parboxrestore{\RaggedRight}
\makeatother

\tikzset{
  vecteur/.style={-{Stealth[length=2.4mm]},line width=1pt,coulCourbe},
  vecteurB/.style={-{Stealth[length=2.4mm]},line width=1pt,coulCourbeB},
  vecteurC/.style={-{Stealth[length=2.4mm]},line width=1pt,coulCourbeC},
  fig/.style={line width=0.8pt,black!75},
}

\begin{document}

\titrefiche{Seconde --- Vecteurs}{Fiche d'exercices du chapitre}

{\small\itshape Les figures se font au crayon et à la règle. On reproduit le
quadrillage quand il y en a un, et on vérifie chaque vecteur en comptant les
carreaux.\par}
\vspace{0.12cm}

\begin{multicols}{2}\raggedcolumns

\partiefiche{Caractéristiques d'un vecteur}

\exo[Égaux ou non]
La figure ci-dessous est tracée sur un quadrillage.
\begin{center}
\begin{tikzpicture}[scale=0.5]
  \quadrillage{-0.4}{-0.4}{8.4}{2.4}
  \draw[fig] (0,0)--(6,0) (2,2)--(8,2) (0,0)--(2,2) (3,0)--(5,2) (6,0)--(8,2);
  \pt{0,0}{below}{K} \pt{3,0}{below}{L} \pt{6,0}{below}{M}
  \pt{2,2}{above}{N} \pt{5,2}{above}{P} \pt{8,2}{above}{Q}
\end{tikzpicture}
\end{center}
Compléter par $=$ ou $\neq$ :
\begin{enumerate}
  \item[a)] $\vect{KL}\;\ldots\;\vect{NP}$ \quad b) $\vect{KN}\;\ldots\;\vect{LP}$
  \item[c)] $\vect{KP}\;\ldots\;\vect{LQ}$ \quad d) $\vect{NL}\;\ldots\;\vect{KN}$
  \item[e)] $\vect{KM}\;\ldots\;\vect{NQ}$ \quad f) $\vect{PQ}\;\ldots\;\vect{PN}$
\end{enumerate}

\exo[Lire sur la figure]
On utilise la figure de l'exercice 1.
\begin{enumerate}
  \item Citer un vecteur égal à $\vect{KN}$, puis un vecteur égal à $\vect{LM}$.
  \item Citer un vecteur opposé à $\vect{KN}$, puis un vecteur opposé à $\vect{KL}$.
  \item Citer deux vecteurs :
  \begin{itemize}
    \item de même direction et de même sens, mais pas égaux ;
    \item de même norme, mais pas égaux ;
    \item de même direction et de sens contraires, mais pas opposés.
  \end{itemize}
\end{enumerate}

\exo[Construire un représentant]
Tracer un vecteur $\vec{u}$ de norme $3$~cm. Construire ensuite un représentant du
vecteur $\vec{v}$ dans chaque cas :
\begin{enumerate}
  \item $\vec{v}$ a la même direction et le même sens que $\vec{u}$, et pour norme
        $6$~cm ;
  \item $\vec{v}$ a la même direction et la même norme que $\vec{u}$, et le sens
        contraire ;
  \item $\vec{v}$ a la même direction que $\vec{u}$, le sens contraire, et pour norme
        $1{,}5$~cm ;
  \item $\vec{v}$ a la même norme que $\vec{u}$ et une autre direction.
\end{enumerate}

\exo[Vrai ou faux]
$EFGH$ et $HGKL$ sont deux parallélogrammes. Dire si chaque affirmation est vraie ou
fausse, en justifiant.
\begin{enumerate}
  \item L'image de $H$ par la translation de vecteur $\vect{GF}$ est le point $E$.
  \item Les vecteurs $\vect{KG}$ et $\vect{HL}$ ont la même direction.
  \item $\vect{KG}=\vect{HL}$.
  \item $\vect{EF}=\vect{LK}$.
\end{enumerate}

\exo[Dans un parallélogramme]
$RSTU$ est un parallélogramme de centre $O$. Faire une figure.
\begin{enumerate}
  \item Donner un vecteur égal à :
        \choix{$\vect{RS}$}{$\vect{TU}$}{$\vect{RU}$}
  \item Compléter par le point qui convient :
        \choix{$\vect{RO}=\vect{O\ldots}$}{$\vect{UO}=\vect{O\ldots}$}{$\vect{SO}=\vect{O\ldots}$}
\end{enumerate}

\partiefiche{Translation}

\exo[Représentants d'origine donnée]
Tracer un triangle $DEF$ et un point $P$. Tracer le représentant d'origine $P$ de
chacun des vecteurs $\vect{DE}$, $\vect{EF}$ et $\vect{FD}$.

\exo[Quatre points à placer]
$ABC$ est un triangle. Placer les points $M$, $N$, $P$ et $Q$ tels que
\[ \vect{AC}=\vect{BM}=\vect{NB}, \qquad \vect{PC}=\vect{AB}, \qquad \vect{AQ}=\vect{CB}. \]

\exo[Origine ou extrémité]
On donne un vecteur non nul $\vec{u}$ et trois points $E$, $F$, $G$. Tracer le
représentant de $\vec{u}$ d'origine $E$, puis celui d'\emph{extrémité} $F$, puis celui
d'\emph{extrémité} $G$.

\exo[Image d'un triangle]
Tracer un triangle $ABC$ tel que $AB=5$~cm, $BC=4$~cm et $AC=3$~cm. Construire les
images $A'$, $B'$, $C'$ de ses sommets par la translation de vecteur $\vect{BC}$. Que
dire des triangles $ABC$ et $A'B'C'$ ?

\exo[Neuf points]
Neuf points sont placés sur un quadrillage.
\begin{center}
\begin{tikzpicture}[scale=0.45]
  \quadrillage{-0.4}{-0.4}{4.4}{4.4}
  \pt{0,4}{above left}{A} \pt{2,4}{above}{B} \pt{4,4}{above right}{C}
  \pt{0,2}{left}{D} \pt{2,2}{above right}{E} \pt{4,2}{right}{F}
  \pt{0,0}{below left}{G} \pt{2,0}{below}{H} \pt{4,0}{below right}{I}
\end{tikzpicture}
\end{center}
\begin{enumerate}
  \item Citer tous les vecteurs formés par deux de ces points et égaux à $\vect{AB}$.
  \item Citer tous ceux qui sont égaux à $\vect{AE}$.
  \item Citer deux vecteurs opposés à $\vect{AB}$.
\end{enumerate}

\partiefiche{Parallélogramme et égalité de vecteurs}

\exo[Deux parallélogrammes]
$EFGH$ et $EFIJ$ sont deux parallélogrammes. Faire une figure.
\begin{enumerate}
  \item Donner, en justifiant, deux vecteurs égaux à $\vect{EF}$.
  \item Démontrer que $HGIJ$ est un parallélogramme.
\end{enumerate}

\exo[Autour d'un triangle]
Tracer un triangle $DEF$.
\begin{enumerate}
  \item Construire les points $M$ et $N$ tels que $\vect{FM}=\vect{DE}$ et
        $\vect{NF}=\vect{DE}$.
  \item Construire le point $K$, image de $D$ par la translation de vecteur $\vect{FE}$.
  \item Quel point de la figure a pour image $E$ par la translation de vecteur
        $\vect{FD}$ ?
  \item Écrire trois autres égalités de vecteurs lues sur la figure.
\end{enumerate}

\exo[Un milieu à démontrer]
$RSTU$ est un parallélogramme.
\begin{enumerate}
  \item Faire une figure.
  \item Construire le point $V$ tel que $RSUV$ soit un parallélogramme.
  \item Démontrer que $U$ est le milieu de $[TV]$.
\end{enumerate}

\exo[Vrai ou faux]
Dire si chaque affirmation est vraie ou fausse, en justifiant.
\begin{enumerate}
  \item Si $F$ est l'image de $C$ par la translation de vecteur $\vect{DE}$, alors
        $\vect{FC}=\vect{DE}$.
  \item Si $\vect{KL}=\vect{NM}$, alors $KLMN$ est un parallélogramme.
  \item Si $ABCD$ est un parallélogramme, la translation de vecteur $\vect{AD}$
        transforme $B$ en $C$.
\end{enumerate}

\partiefiche{Somme de deux vecteurs}

\exo[Trois points à construire]
Tracer un triangle $ABC$, puis construire les points $D$, $E$ et $F$ tels que
\[ \vect{BD}=\vect{BC}+\vect{CA}, \quad \vect{BE}=\vect{BA}+\vect{BC}, \quad
   \vect{CF}=\vect{CB}+\vect{AB}. \]

\exo[Somme et différence sur un quadrillage]
Reproduire chaque figure, puis placer le point $M$ tel que $\vect{AM}=\vec{u}+\vec{v}$
et le point $N$ tel que $\vect{AN}=\vec{u}-\vec{v}$.
\begin{center}
\begin{tikzpicture}[scale=0.45]
  \quadrillage{-0.4}{-0.4}{5.4}{4.4}
  \pt{1,1}{below left}{A}
  \draw[vecteur] (1,1)--(4,2) node[midway,below right,font=\small]{$\vec{u}$};
  \draw[vecteurB] (1,1)--(2,3) node[midway,left,font=\small]{$\vec{v}$};
\end{tikzpicture}
\hspace{6pt}
\begin{tikzpicture}[scale=0.45]
  \quadrillage{-0.4}{-0.4}{5.4}{4.4}
  \pt{1,2}{left}{A}
  \draw[vecteur] (1,2)--(4,2) node[midway,above,font=\small]{$\vec{u}$};
  \draw[vecteurB] (1,2)--(1,0) node[midway,left,font=\small]{$\vec{v}$};
\end{tikzpicture}
\end{center}

\exo[Origine imposée]
On pose $\vec{w}=\vec{u}+\vec{v}$. Reproduire chaque figure et tracer le représentant de
$\vec{w}$ d'origine $T$.
\begin{center}
\begin{tikzpicture}[scale=0.45]
  \quadrillage{-0.4}{-0.4}{5.4}{4.4}
  \pt{1,1}{below left}{T}
  \draw[vecteur]  (0,2)--(1,4) node[midway,left,font=\small]{$\vec{u}$};
  \draw[vecteurB] (3,4)--(5,3) node[midway,above right,font=\small]{$\vec{v}$};
\end{tikzpicture}
\hspace{6pt}
\begin{tikzpicture}[scale=0.45]
  \quadrillage{-0.4}{-0.4}{5.4}{4.4}
  \pt{3,4}{above}{T}
  \draw[vecteur]  (0,4)--(2,3) node[midway,below left,font=\small]{$\vec{u}$};
  \draw[vecteurB] (2,2)--(1,0) node[midway,left,font=\small]{$\vec{v}$};
\end{tikzpicture}
\end{center}

\exo[Quand c'est l'origine qui manque]
Ici, le point à placer est l'\textbf{origine} du vecteur, et non son extrémité.
Reproduire chaque figure, puis placer le point $N$ tel que :
\begin{enumerate}
  \item $\vect{NA}=\vec{u}+\vec{v}$ \quad (figure de gauche) ;
  \item $\vect{NA}=\vec{u}-\vec{v}$ \quad (figure de droite).
\end{enumerate}
\begin{center}
\begin{tikzpicture}[scale=0.45]
  \quadrillage{-0.4}{-0.4}{5.4}{4.4}
  \pt{4,2}{below right}{A}
  \draw[vecteur]  (0,1)--(1,3) node[midway,left,font=\small]{$\vec{u}$};
  \draw[vecteurB] (2,4)--(4,3) node[midway,above right,font=\small]{$\vec{v}$};
\end{tikzpicture}
\hspace{6pt}
\begin{tikzpicture}[scale=0.45]
  \quadrillage{-0.4}{-0.4}{5.4}{4.4}
  \pt{2,2}{below left}{A}
  \draw[vecteur]  (0,3)--(2,4) node[midway,above left,font=\small]{$\vec{u}$};
  \draw[vecteurB] (4,0)--(5,2) node[midway,right,font=\small]{$\vec{v}$};
\end{tikzpicture}
\end{center}

\exo[Avec des longueurs]
$ABC$ est un triangle tel que $AB=6$~cm, $AC=5$~cm et $BC=4$~cm. Faire une figure en
vraie grandeur.
\begin{enumerate}
  \item Construire le point $D$ tel que $\vect{AD}=\vect{AB}+\vect{AC}$.
  \item Construire le point $E$ tel que $\vect{BE}=\vect{BC}-\vect{BA}$.
  \item Quelle est la nature du quadrilatère $ABDC$ ? Justifier.
\end{enumerate}

\exo[Trois vecteurs]
Reproduire la figure, puis construire :
\begin{enumerate}
  \item les représentants d'origine $A$ de $\vec{u}+\vec{v}$ et de $\vec{u}+\vec{w}$ ;
  \item le représentant d'origine $B$ de $\vec{v}+\vec{w}$.
\end{enumerate}
\begin{center}
\begin{tikzpicture}[scale=0.45]
  \quadrillage{-0.4}{-0.4}{8.4}{4.4}
  \pt{1,1}{below left}{A} \pt{5,0}{below}{B}
  \draw[vecteur]  (2,3)--(3,4) node[midway,left,font=\small]{$\vec{u}$};
  \draw[vecteurB] (4,4)--(6,3) node[midway,above right,font=\small]{$\vec{v}$};
  \draw[vecteurC] (8,1)--(7,3) node[midway,right,font=\small]{$\vec{w}$};
\end{tikzpicture}
\end{center}

\partiefiche{Milieu d'un segment}

\exo[Une différence, un parallélogramme]
$ABC$ est un triangle.
\begin{enumerate}
  \item Construire le point $D$ tel que $\vect{AD}=\vect{AC}-\vect{AB}$. Quelle est la
        nature du quadrilatère $ADCB$ ? Que dire des segments $[AC]$ et $[DB]$ ?
  \item Placer le point $E$ tel que $\vect{CE}=\vect{BA}+\vect{CB}$. Que constate-t-on ?
\end{enumerate}

\exo[Trois points à construire]
$ABC$ est un triangle. Construire le point $M$ tel que :
\begin{enumerate}
  \item $\vect{AM}=\vect{AC}-\vect{AB}$ ;
  \item $\vect{CM}=\vect{CB}-\vect{BA}$ ;
  \item $\vect{BM}=\vect{AC}+\vect{BA}-\vect{CB}$.
\end{enumerate}

\exo[Symétrie centrale]
$E$, $F$ et $G$ sont trois points. $E'$ est le symétrique de $E$ par rapport à $G$, et
$F'$ celui de $F$ par rapport à $G$.
\begin{enumerate}
  \item Démontrer que $FEF'E'$ est un parallélogramme.
  \item \begin{enumerate}
          \item Justifier que $\vect{E'F'}=\vect{FE}$.
          \item Que peut-on en déduire pour les longueurs $E'F'$ et $EF$ ?
          \item Quelle propriété de la symétrie centrale vient-on de démontrer ?
        \end{enumerate}
\end{enumerate}

\partiefiche{Multiplier un vecteur par un réel}

\exo[Multiples d'un vecteur]
Reproduire la figure, puis représenter, à partir du point $O$, les vecteurs
$2\vec{a}$, $3\vec{a}$, $-\vec{a}$ et $-2\vec{a}$.
\begin{center}
\begin{tikzpicture}[scale=0.45]
  \quadrillage{-0.4}{-0.4}{7.4}{5.4}
  \pt{3,3}{left}{O}
  \draw[vecteur] (3,3)--(4,2) node[midway,above right,font=\small]{$\vec{a}$};
\end{tikzpicture}
\end{center}

\exo[Avec des coefficients]
$ABC$ est un triangle. Construire les points $M$, $N$ et $P$ tels que
\[ \vect{AM}=3\vect{AC}, \qquad \vect{CN}=\tfrac12\,\vect{CB}, \qquad
   \vect{BP}=-\tfrac12\,\vect{BA}. \]

\exo[Sur un quadrillage]
Quatre points $A$, $B$, $C$ et $D$ sont placés sur un quadrillage.
\begin{center}
\begin{tikzpicture}[scale=0.45]
  \quadrillage{-0.4}{-0.4}{8.4}{5.4}
  \pt{0,1}{below left}{A} \pt{4,1}{below}{B}
  \pt{6,3}{right}{C} \pt{1,4}{above left}{D}
\end{tikzpicture}
\end{center}
Reproduire la figure, puis tracer :
\begin{enumerate}
  \item le représentant d'origine $A$ de $2\vect{BC}$ ;
  \item le représentant d'origine $B$ de $0{,}5\,\vect{AC}$ ;
  \item le représentant d'origine $D$ de $-\tfrac13\,\vect{BD}$ ;
  \item le représentant d'origine $C$ de $\tfrac13\,\vect{BD}$ ;
  \item le représentant d'origine $D$ de $\vect{AB}-\vect{BC}$.
\end{enumerate}
\emph{Chaque extrémité tombe sur un n\oe ud du quadrillage.}

\partiefiche{Relation de Chasles}

\exo[Un seul vecteur, si c'est possible]
Écrire avec un seul vecteur lorsque c'est possible ; sinon, expliquer pourquoi.
\begin{enumerate}
  \item[a)] $\vect{BE}+\vect{EG}$ \quad b) $\vect{DA}+\vect{DC}$
  \item[c)] $\vect{TU}+\vect{KT}$ \quad d) $\vect{FH}+\vect{HS}$
  \item[e)] $\vect{PQ}+\vect{QP}$ \quad f) $\vect{AB}+\vect{CE}$
\end{enumerate}

\exo[Simplifier]
$ABC$ est un triangle et $M$ un point. Simplifier à l'aide de la relation de Chasles :
\[ \vec{u}=\vect{BC}+\vect{CA}+\vect{AB} \qquad \vec{v}=\vect{CA}-\vect{CB}+\vect{AC} \]
\[ \vec{w}=\vect{MC}-\vect{MA}+\vect{CB} \]
\emph{Aide} : $-\vect{CB}=\vect{BC}$.

\exo[Compléter]
Compléter par les points qui conviennent.
\begin{enumerate}
  \item $\vect{BC}=\vect{BA}+\vect{A\ldots}$
  \item $\vect{CB}+\vect{\ldots A}=\vect{CA}$
  \item $\vect{A\ldots}+\vect{CA}=\vect{CB}$
  \item $\vect{BC}+\vect{DA}+\vect{CD}=\vect{\ldots}$
\end{enumerate}

\exo[Peut-on réduire ?]
Dire si chaque somme se réduit avec la relation de Chasles, et la réduire si c'est le
cas.
\begin{enumerate}
  \item[a)] $\vect{MN}+\vect{NP}$ \quad b) $\vect{MN}+\vect{MP}$
  \item[c)] $\vect{RO}+\vect{OS}$ \quad d) $\vect{RO}+\vect{SR}$
  \item[e)] $\vect{GH}+\vect{HG}$ \quad f) $\vect{KB}+\vect{AK}$
\end{enumerate}

\exo[Une égalité à démontrer]
Démontrer que, pour tous points $O$, $A$, $B$ et $C$ :
\[ \vect{OB}-\vect{OA}+\vect{BC}=\vect{AC}. \]

\exo[Exprimer avec deux vecteurs]
$ABC$ est un triangle.
\begin{enumerate}
  \item Exprimer $\vec{u}$ en fonction de $\vect{AB}$ et $\vect{AC}$ :
        \par\hspace*{1em}a) $\vec{u}=\vect{CB}$
        \quad b) $\vec{u}=3\vect{CB}+\vect{BA}$
  \item Exprimer $\vec{v}$ en fonction de $\vect{CA}$ et $\vect{BC}$ :
        \par\hspace*{1em}a) $\vec{v}=\vect{BA}+\vect{CA}$
        \par\hspace*{1em}b) $\vec{v}=\vect{AB}+2\vect{BC}+\vect{AC}$
\end{enumerate}

\exo[Différences sur un quadrillage]
Construire, à partir du point $A$, un représentant de chaque vecteur demandé.
\emph{On peut prolonger le quadrillage sur la feuille.}
\begin{center}
\begin{tikzpicture}[scale=0.45]
  \quadrillage{-0.4}{-0.4}{6.4}{3.4}
  \pt{0,1}{left}{A} \pt{4,0}{below}{B}
  \pt{5,3}{above right}{C} \pt{1,3}{above}{D}
\end{tikzpicture}
\end{center}
\begin{enumerate}
  \item[a)] $\vect{AC}-\vect{AB}$ \quad b) $\vect{AD}-\vect{BD}$
  \item[c)] $\vect{AB}-\vect{AD}$ \quad d) $\vect{AC}-\vect{BD}$
\end{enumerate}

\end{multicols}

\end{document}
